% GENERATED FILE. Edit canonical lessons, then run npm run generate:books.
% canonical-source-hash: fnv1a64:d40de39d7a9852ba
% canonical-lessons: TE-C195-adhikari, TE-C195-gumasta, TE-C195-menejaru, TE-C195-sahodyogi, TE-C195-samavesam

\chapter{Officer, Clerk, Manager, Colleague, Meeting}
\label{ch:te-officer-clerk-manager-colleague-meeting}
\hlchaptermodality{195}

\begin{chapteropening}
\textbf{By the end of this chapter:} \emph{I can say an officer, a clerk, a manager, a colleague and a meeting.}

Complete the last lesson of chapter 195: I can say an officer, a clerk, a manager, a colleague and a meeting.
\end{chapteropening}

\section[\texorpdfstring{\te{అధికారి} (adhikāri)}{అధికారి (adhikāri)}]{\te{అధికారి} (adhikāri) — an officer, an official}

\label{lesson:TE-C195-adhikari}



\begin{quote}
Before the new one: say the Telugu for a trader, then the Telugu for a coin.
\end{quote}

\subsection*{You'll want to know: \te{అధికారి}}
\textbf{\te{అధికారి}} — \emph{adhikāri} — \textquotedblleft{}an officer, an official\textquotedblright{}.

\begin{grammarlens}[title={What you've built}]
One more word for this chapter.
\end{grammarlens}

\subsection*{Your turn}
\begin{itemize}
\raggedright
  \item \emph{Say it:} \emph{adhikāri}
  \item \emph{Say it:} \emph{adhikāri}, once more
  \item \emph{Say it:} say \emph{adhikāri}
  \item \emph{Recall:} say the Telugu for a trader, then the Telugu for a coin, then say \emph{adhikāri} again
\end{itemize}

\begin{tcolorbox}[breakable,colback=teal!4,colframe=teal!35!black,title={Before you move on}]
What does \te{అధికారి} mean? (\textquotedblleft{}An officer, an official\textquotedblright{}.) Say it once more.
\end{tcolorbox}
\section[\texorpdfstring{\te{గుమాస్తా} (gumāstā)}{గుమాస్తా (gumāstā)}]{\te{గుమాస్తా} (gumāstā) — a clerk}

\label{lesson:TE-C195-gumasta}



\begin{quote}
Before the new one: say the Telugu for a coin, then the Telugu for an officer.
\end{quote}

\subsection*{You'll want to know: \te{గుమాస్తా}}
\textbf{\te{గుమాస్తా}} — \emph{gumāstā} — \textquotedblleft{}a clerk\textquotedblright{}.

\begin{grammarlens}[title={What you've built}]
One more word for this chapter.
\end{grammarlens}

\subsection*{Your turn}
\begin{itemize}
\raggedright
  \item \emph{Say it:} \emph{gumāstā}
  \item \emph{Say it:} \emph{gumāstā}, once more
  \item \emph{Say it:} say \emph{gumāstā}
  \item \emph{Recall:} say the Telugu for a coin, then the Telugu for an officer, then say \emph{gumāstā} again
\end{itemize}

\begin{tcolorbox}[breakable,colback=teal!4,colframe=teal!35!black,title={Before you move on}]
What does \te{గుమాస్తా} mean? (\textquotedblleft{}A clerk\textquotedblright{}.) Say it once more.
\end{tcolorbox}
\section[\texorpdfstring{\te{మేనేజరు} (mēnējaru)}{మేనేజరు (mēnējaru)}]{\te{మేనేజరు} (mēnējaru) — a manager}

\label{lesson:TE-C195-menejaru}



\begin{quote}
Before the new one: say the Telugu for an officer, then the Telugu for a clerk.
\end{quote}

\subsection*{You'll want to know: \te{మేనేజరు}}
\textbf{\te{మేనేజరు}} — \emph{mēnējaru} — \textquotedblleft{}a manager\textquotedblright{}.

\begin{grammarlens}[title={What you've built}]
One more word for this chapter.
\end{grammarlens}

\subsection*{Your turn}
\begin{itemize}
\raggedright
  \item \emph{Say it:} \emph{mēnējaru}
  \item \emph{Say it:} \emph{mēnējaru}, once more
  \item \emph{Say it:} say \emph{mēnējaru}
  \item \emph{Recall:} say the Telugu for an officer, then the Telugu for a clerk, then say \emph{mēnējaru} again
\end{itemize}

\begin{tcolorbox}[breakable,colback=teal!4,colframe=teal!35!black,title={Before you move on}]
What does \te{మేనేజరు} mean? (\textquotedblleft{}A manager\textquotedblright{}.) Say it once more.
\end{tcolorbox}
\section[\texorpdfstring{\te{సహోద్యోగి} (sahōdyōgi)}{సహోద్యోగి (sahōdyōgi)}]{\te{సహోద్యోగి} (sahōdyōgi) — a colleague}

\label{lesson:TE-C195-sahodyogi}



\begin{quote}
Before the new one: say the Telugu for a clerk, then the Telugu for a manager.
\end{quote}

\subsection*{You'll want to know: \te{సహోద్యోగి}}
\textbf{\te{సహోద్యోగి}} — \emph{sahōdyōgi} — \textquotedblleft{}a colleague\textquotedblright{}.

\begin{grammarlens}[title={What you've built}]
One more word for this chapter.
\end{grammarlens}

\subsection*{Your turn}
\begin{itemize}
\raggedright
  \item \emph{Say it:} \emph{sahōdyōgi}
  \item \emph{Say it:} \emph{sahōdyōgi}, once more
  \item \emph{Say it:} say \emph{sahōdyōgi}
  \item \emph{Recall:} say the Telugu for a clerk, then the Telugu for a manager, then say \emph{sahōdyōgi} again
\end{itemize}

\begin{tcolorbox}[breakable,colback=teal!4,colframe=teal!35!black,title={Before you move on}]
What does \te{సహోద్యోగి} mean? (\textquotedblleft{}A colleague\textquotedblright{}.) Say it once more.
\end{tcolorbox}
\section[\texorpdfstring{\te{సమావేశం} (samāvēśaṁ)}{సమావేశం (samāvēśaṁ)}]{\te{సమావేశం} (samāvēśaṁ) — a meeting}

\label{lesson:TE-C195-samavesam}



\begin{quote}
Before the new one: say the Telugu for a manager, then the Telugu for a colleague.
\end{quote}

\subsection*{You'll want to know: \te{సమావేశం}}
\textbf{\te{సమావేశం}} — \emph{samāvēśaṁ} — \textquotedblleft{}a meeting\textquotedblright{}.

\begin{grammarlens}[title={What you've built}]
One more word for this chapter.
\end{grammarlens}

\subsection*{Your turn}
\begin{itemize}
\raggedright
  \item \emph{Say it:} \emph{samāvēśaṁ}
  \item \emph{Say it:} \emph{samāvēśaṁ}, once more
  \item \emph{Say it:} say \emph{samāvēśaṁ}
  \item \emph{Recall:} say the Telugu for a manager, then the Telugu for a colleague, then say \emph{samāvēśaṁ} again
\end{itemize}

\begin{tcolorbox}[breakable,colback=teal!4,colframe=teal!35!black,title={Before you move on}]
What does \te{సమావేశం} mean? (\textquotedblleft{}A meeting\textquotedblright{}.) Say it once more.
\end{tcolorbox}
